\clearpage
\VerdierPhysicalPageAnchor{0119}%
\begingroup
\small
\setlength{\abovedisplayskip}{4pt plus 2pt minus 2pt}%
\setlength{\belowdisplayskip}{4pt plus 2pt minus 2pt}%

\VerdierSemanticAnchor{ch2-1-2-8-splitting-complete}
$fgf=f$ 且 $gfg=g$。一个态射称为\emph{可分裂的}，如果它具有一个分裂。
可分解态射都是可分裂的；反之未必成立（例如：非主理想 Dedekind 环上
自由模的范畴）。

\VerdierSemanticAnchor{ch2-1-2-8-decomposable-category}
一个加性范畴称为\emph{可分解的}，如果其中每个态射都可分解。
可分解范畴是阿贝尔范畴。

\VerdierSemanticAnchor{ch2-1-2-8-triangle-decomposition}
设 $X\xrightarrow{f}Y$ 为三角范畴 $\mathcal D$ 中的态射，并设

\VerdierSemanticAnchor{ch2-1-2-8-decomposition-diagram-repeat}
\begin{equation*}
  \begin{tikzcd}[ampersand replacement=\&,column sep=5.2em,row sep=3.4em]
    X \arrow[r,"f"] \arrow[d,"\wr"']
      \& Y \arrow[d,"\wr"] \\
    Z\oplus W
      \arrow[r,"{\left(\begin{smallmatrix}0&\operatorname{id}_W\\0&0\end{smallmatrix}\right)}"]
      \& W\oplus Z'
  \end{tikzcd}\,.
\end{equation*}
为 $f$ 的一个分解。每个包含 $f$ 的特异三角都同构于三角

\VerdierSemanticAnchor{ch2-1-2-8-decomposed-triangle}
\begin{equation*}
  \begin{tikzcd}[ampersand replacement=\&,column sep=2.0em]
    Z\oplus W
      \arrow[r,"{\left(\begin{smallmatrix}0&\operatorname{id}_W\\0&0\end{smallmatrix}\right)}"]
      \& W\oplus Z'
      \arrow[r,"{\left(\begin{smallmatrix}0&\operatorname{id}_{Z'}\\0&0\end{smallmatrix}\right)}"]
      \& Z'\oplus Z[1]
      \arrow[r,"{\left(\begin{smallmatrix}0&\operatorname{id}_{Z[1]}\\0&0\end{smallmatrix}\right)}"]
      \& Z[1]\oplus W[1]
  \end{tikzcd}\,.
\end{equation*}

\VerdierSemanticAnchor{ch2-1-2-8-decomposed-triangle-derivation}
这立即由 (1.2.4)、(1.2.5) 以及公理 TRI、TRII 得到。

\VerdierSemanticAnchor{ch2-1-2-9}
\medskip
\noindent\textbf{命题 1.2.9.}\quad
\emph{设 $\mathcal D$ 为三角范畴，$f\colon X\longrightarrow Y$ 为
$\mathcal D$ 中的态射。考虑 $f$ 的下列性质：}

\VerdierSemanticAnchor{ch2-1-2-9-properties}
\begin{enumerate}
\renewcommand{\labelenumi}{\roman{enumi})}
\item \emph{态射 $f$ 有核。}
\item \emph{态射 $f$ 有余核。}
\item \emph{态射 $f$ 可分解。}
\item \emph{态射 $f$ 可分裂。}
\end{enumerate}

\noindent\emph{则有蕴含关系：}

\VerdierSemanticAnchor{ch2-1-2-9-implications}
\[
  \mathrm{i)}\Longleftrightarrow\mathrm{ii)}
  \Longleftrightarrow\mathrm{iii)}\Longrightarrow\mathrm{iv)}\,.
\]

\VerdierSemanticAnchor{ch2-1-2-9-countable-equivalence}
\emph{如果可数积或可数直和在 $\mathcal D$ 中可表示，则性质
\textnormal{i)}、\textnormal{ii)}、\textnormal{iii)}、\textnormal{iv)}
等价。}

\VerdierSemanticAnchor{ch2-1-2-9-proof-opening}
首先，显然有蕴含
$\mathrm{iii)}\Longrightarrow\mathrm{i)}$、
$\mathrm{iii)}\Longrightarrow\mathrm{ii)}$、
$\mathrm{iii)}\Longrightarrow\mathrm{iv)}$。蕴含
$\mathrm{i)}\Longrightarrow\mathrm{iii)}$ 等价于蕴含
\endgroup
